\chapter{Conjuntos Abstractos y Espacios}
\label{conjuntos}

Originalmente, Lawvere propuso los siguientes $8$ axiomas para caracterizar a una categoría
de conjuntos abstractos.

\begin{axiom}[Raíces finitas]
Existen todas las raíces (límites y colimites) finitas. Explícitamente:
existen un objeto terminal $1$ y un objeto inicial $0$, para cada par de objetos
$A,B$ existen el producto $A\times B$ y la suma (coproducto) $A+B$, y para cada
par de flechas paralelas $f,g\colon A\to B$ existen su igualador y su coigualador.
\end{axiom}

\begin{axiom}[Exponenciación]
Para cada par de objetos $A,B$ existe el objeto exponencial $B^{A}$.
\end{axiom}

\begin{axiom}[Objeto de Dedekind--Peano]
Existe un objeto $(N,o,s)$ de números naturales
\end{axiom}

\begin{axiom}[$1$ es un separador]
Si $f,g\colon A\to B$ son flechas distintas, entonces existe un elemento
$x\colon 1\to A$ tal que $fx\neq gx$.
\end{axiom}

\begin{axiom}[Axioma de elección]
Si $f\colon A\to B$ es una flecha y $A$ tiene al menos un elemento, entonces
existe una flecha $g\colon B\to A$ tal que $fgf=f$.
\end{axiom}

\begin{axiom}[Los objetos no iniciales tienen elementos]
Si $A$ no es un objeto inicial, entonces existe al menos un elemento
$x\colon 1\to A$.
\end{axiom}

\begin{axiom}[Elementos de una suma]
Sean $i_A\colon A\to A+B$ e $i_B\colon B\to A+B$ las inyecciones canónicas.
Todo elemento $x\colon 1\to A+B$ factoriza a través de $i_A$ o a través de
$i_B$.
\end{axiom}

\begin{axiom}[Existencia de un objeto con más de un elemento]
Existe un objeto que tiene más de un elemento. Equivalentemente,
$1$ y $0$ no son isomorfos.
\end{axiom}

Con el desarrollo de la teoría de topos se demostró que la lista se puede simplificar a pedir
que la categoría sea un topos bien punteado, con objeto de números naturales y axioma de elección.